\section{Kognitive Kontrolle / Exekutive}
\slide{Begrifflichkeiten}{
    \begin{itemize}
        \item \b{Performanz:} Die messbare Leistung bei einer Aufgabe, meist operationalisiert durch Reaktionszeit, Fehlerrate oder Präzision.
        \item \b{Interferenz:} Verschlechterung der Performanz bei Doppelaufgaben, wenn beide Aufgaben dieselbe begrenzte Ressource beanspruchen .
        \item \b{Inhibition:} Das aktive Unterdrücken eines nicht aufgabenrelevanten Stimulus oder einer automatisierten Antwort .
        \item \b{Automatische Prozesse:} Abläufe, die schnell und ohne bewusste Aufmerksamkeit durch abrufbare Schemata gesteuert werden (z. B. Lesen).
        \item \b{Kontrollierte Prozesse:} Abläufe, die Aufmerksamkeit erfordern, flexibel sind und bei neuen Situationen ohne vorgefertigte Schemata greifen .
        \item \b{Metakognition:} Die Überwachung und Reflexion der eigenen kognitiven Prozesse.
    \end{itemize}
}
\slide{Ressourcentheorien}{
    \begin{itemize}
        \item \b{Ressourcentheorie} (Kahneman): Es gibt eine zentrale kognitive Ressource mit begrenzter Kapazität, die bei Multi-Tasking aufgeteilt werden muss.
        \item \b{Theorie multipler Ressourcen} (Navon, Gopher, Wickens): Es gibt mehrere domänen-spezifische Ressourcen (z. B. visuell vs. auditiv, verbal vs. motorisch) mit unabhängigen Kapazitäten.
        \begin{itemize}
            \item Interferenz tritt nur auf, wenn zwei Aufgaben dieselbe spezifische Ressource nutzen.
        \end{itemize}
    \end{itemize}
}
\slide{Doppelaufgaben-Paradigma (Dual Task)}{
    \begin{itemize}
        \item \b{Rekognitionsexperiment} (Formen \& Zählen):
        \begin{itemize}
            \item Aufgabe: Wiedererkennen abstrakter Formen bei gleichzeitiger Zweitaufgabe (Rückwärtszählen).
            \item Interpretation: Findet keine Interferenz statt, wurden die Formen visuell enkodiert (unterschiedliche Ressourcen). Tritt Interferenz auf, wurden die Formen verbal enkodiert (gleiche Ressource wie das Zählen).
        \end{itemize}
        \item \b{Spelke, Hirst \& Neisser} (Lesen \& Schreiben):
        \begin{itemize}
            \item Experiment: Gleichzeitiges Lesen von Geschichten (Verständnis) und Schreiben diktierter Wörter über mehrere Wochen.
            \item Ergebnis: Performanz-Steigerung durch Übung.
            \item Interpretation: Die Ressourcennutzung wurde effizienter, oder eine Aufgabe wurde automatisiert, wodurch weniger kognitiver Aufwand nötig war.
        \end{itemize}
    \end{itemize}
}
\slide{PRP-Paradigma (Psychologische Refraktär-Periode)}{
    \begin{itemize}
        \item Funktionsweise: Zwei Aufgaben (z. B. Tonhöhe unterscheiden + Form erkennen) werden kurz nacheinander präsentiert. Das Zeitintervall zwischen den Reizen (SOA) wird variiert.
        \item Hypothesen zur Schwierigkeit:
        \begin{enumerate}
            \item Erhöhte Schwierigkeit verlängert die kognitive Verarbeitung (Phase B).
            \item Erhöhte Schwierigkeit verlängert die Wahrnehmung/Enkodierung (Phase A).
        \end{enumerate}
        \item Ergebnisse: Hypothese 1 wurde bestätigt. Die Reaktionszeit der zweiten Aufgabe verlängert sich bei schwierigen Aufgaben parallel zur Bearbeitungszeit der ersten, besonders bei kurzer SOA.
        \item Flaschenhals: Es existiert ein Engpass in der kognitiven Verarbeitung; diese ist auf einen Reiz beschränkt (serielle Verarbeitung), während Wahrnehmung und Motorik teilweise parallel laufen können.
    \end{itemize}
}
\slide{Task-Switching (Aufgabenwechsel)}{
    \begin{itemize}
        \item Erhöhung der Wechselkosten:
        \begin{itemize}
            \item Unvorhersehbarer Wechsel.
            \item Kurze Zeitintervalle zwischen den Durchgängen.
            \item Überlappung von Aufgaben-, Stimulus- und Antwortmodi.
        \end{itemize}
        \item Reduzierung der Wechselkosten durch mehr Übung und mehr Vorbereitungszeit (aber nie vollständig eliminierbar).
    \end{itemize}
}
\slide{Supervisory Attentional System (SAS) \& Stroop}{
    \begin{itemize}
        \item \b{Stroop-Effekt:} Testet die Inhibition automatischer Prozesse (Lesen) zugunsten kontrollierter Prozesse (Farbe benennen).
        \begin{itemize}
            \item Bei Erwachsenen ist Lesen voll automatisiert, weshalb Interferenz entsteht. Bei Grundschulkindern ist dieser Prozess oft noch nicht automatisiert, weshalb der Effekt schwächer wäre.
        \end{itemize}
        \item \b{SAS-Modell (Norman \& Shallice):}
        \begin{itemize}
            \item Schemata: Erlernte Handlungsabläufe oder kognitive Reaktionen auf spezifische Stimuli.
            \item Voll automatisiert: Stimulus ruft direkt ein spezifisches Schema ab
            \item Partiell automatisiert: "Contention Scheduling" wählt zwischen mehreren möglichen Schemata aus
            \item Kontrolliert (SAS): Das Supervisory Attentional System greift ein, wenn keine Schemata passen oder flexible Reaktionen (wie beim Stroop-Test) nötig sind.
        \end{itemize}
    \end{itemize}
}